% Source pp. 73-79 (Lectures7-12.pdf).
\section{More on Feynman Diagrams. Generating Functional}

Continuing our analysis of the diagrammatic representation of the correlation
functions let us consider the diagram expansion of the four-point function
\[
\bigl\langle\varphi(x_1)\varphi(x_2)\varphi(x_3)\varphi(x_4)\bigr\rangle:=Z^{-1}\int\bigl[D\varphi\bigr]\,\varphi(x_1)\varphi(x_2)\varphi(x_3)\varphi(x_4)\,e^{-\mathcal A[\varphi]}
\]
in the $\varphi^4$ theory. First, by applying the same arguments as before it
is easy to show that in this case (and indeed for any correlation function)
the ``vacuum fragments'' again get canceled by the expansion of $Z$ in the
denominator. However, in this case the disconnected diagrams remain even after
we drop the vacuum fragments. Consider the first few diagrams for the above
four-point function, shown in figure~\ref{fig:fourpt}.
\begin{figure}[ht]
\centering\figfile{fourpt}
\caption{The first few diagrams for the four-point function in the $\varphi^4$
theory. (a)--(c) are the free-theory terms, (d) is a disconnected diagram
with a correction to one of the propagators, and (e), (f) are connected.}
\label{fig:fourpt}
\end{figure}

The first three diagrams here represent the four point function in the free
theory; they consist of the propagators which being connected to the external
points are not connected to each other. These contributions are expressed in a
simple way through the free-field two-point functions $D$. The diagram (d) and
similar diagrams are also disconnected; they represent corrections to those
two-point functions. The diagrams (e) and (f) represent the class of truly
connected diagrams. By the structure of the above diagram series it is natural
to write
\begin{multline*}
\bigl\langle\varphi(x_1)\varphi(x_2)\varphi(x_3)\varphi(x_4)\bigr\rangle
=\bigl\langle\varphi(x_1)\varphi(x_2)\bigr\rangle\bigl\langle\varphi(x_3)\varphi(x_4)\bigr\rangle
+\bigl\langle\varphi(x_1)\varphi(x_3)\bigr\rangle\bigl\langle\varphi(x_2)\varphi(x_4)\bigr\rangle\\
+\bigl\langle\varphi(x_1)\varphi(x_4)\bigr\rangle\bigl\langle\varphi(x_2)\varphi(x_3)\bigr\rangle
+\bigl\langle\varphi(x_1)\varphi(x_2)\varphi(x_3)\varphi(x_4)\bigr\rangle_{\text{conn}}
\end{multline*}
where the first three terms combine the contributions of the disconnected
diagrams; they are expressed in terms of the two-point functions. The connected
diagrams are incorporated in the last term which is called the
\emph{connected} correlation function.

Similar pattern is observed for the higher-order correlation functions. There
are disconnected diagrams which sum up into the lower-order correlation
functions, and besides there are truly connected ones, their sums define the
connected correlation functions.

In fact the disconnected diagrams can be efficiently sorted out by using the so
called Generating Functional.

\subsection*{Generating Functional (PS Sect.~9.2, 9.3)}

Perturbation theory represents any correlation function as the sum of the
Feynman diagrams. There are connected and disconnected diagrams. The vacuum
bubbles can be ignored as they are canceled out by the expansion of the
normalization factor $Z^{-1}$. The remaining disconnected diagrams are
efficiently sorted out by using the method of Generating Functional. Although
the following analysis is done for the case of $\varphi^4$ theory, it is useful
to keep in mind that all the arguments below can be repeated step by step for
the generic interaction.

Let us define the Generating Functional as follows. Introduce the functional
\begin{equation}\label{eq:zetaJ}
\zeta[J]=\sum_{n=0}^\infty\frac1{n!}\int d^4x_1\cdots d^4x_n\,\bigl\langle\varphi(x_1)\cdots\varphi(x_n)\bigr\rangle\,J(x_1)\cdots J(x_n)\,.
\end{equation}
It depends on the auxiliary function $J$, and incorporates all $n$-point
correlation functions. By the definition
\begin{equation}\label{eq:zetader}
\bigl\langle\varphi(x_1)\cdots\varphi(x_n)\bigr\rangle=\frac{\delta}{\delta J(x_1)}\cdots\frac{\delta}{\delta J(x_n)}\,\zeta[J]\Big|_{J=0}\,.
\end{equation}
The generating functional can be written as
\[
\zeta[J]=\Bigl\langle\exp\Bigl(\int J(x)\varphi(x)\,d^4x\Bigr)\Bigr\rangle
\]
and therefore it is expressed as the ratio
\[
\zeta[J]=Z[J]/Z[0]\,;
\]
where
\begin{equation}\label{eq:ZJ}
Z[J]=\int\bigl[D\varphi\bigr]\,e^{\int J(x)\varphi(x)\,d^4x}\,e^{-\mathcal A[\varphi]}
\end{equation}
is a partition function associated with the action
\[
\mathcal A_J=\mathcal A-\int J(x)\varphi(x)\,d^4x\,.
\]
For this reason the function $J(x)$ is usually referred to as an ``external
source''.

The source-dependent action above is a particular case of generic theory with
$x$-dependent coupling $\lambda_1(x)=-J(x)$,
\[
\mathcal A_I=\int d^4x\,\Bigl[-J(x)\varphi(x)+\frac{\lambda}{4!}\,\varphi^4(x)\Bigr]\,.
\]
Therefore its partition function $Z[J]$ (more precisely $Z[J]/Z_0$) is given by
a sum of all vacuum diagrams which now contain two kinds of vertices,
\[
\dia{\draw[prop] (-1.6,0) -- (0,0) node[vtx]{};}=\int d^4x\,J(x)\dots\,,\qquad\qquad
\dia{\pic{xvertex};}=-\lambda\int d^4x\dots\,.
\]
By taking the variational derivatives with respect to $J(x)$ we just convert
the one-leg vertices above into the external legs; this is the meaning of the
formula \eqref{eq:zetader} above. Now, the sum of all the vacuum diagrams
exponentiates in terms of the sum of connected vacuum
diagrams.\footnote{Formal proof of this statement can be found e.g.\ in
``Quantum Field Theory and Critical Phenomena'' by J.~Zinn-Justin, Oxford
Univ.\ Press 1996, \S6.1.} Defining
\begin{equation}\label{eq:Wdef}
\zeta[J]=\exp W[J]
\end{equation}
we conclude that the expansion of this new functional $W[J]$ in terms of $J$
will generate the \emph{connected} correlation functions, namely
\[
W[J]=\sum_{n=0}^\infty\frac1{n!}\int d^4x_1\cdots d^4x_n\,W^{(n)}(x_1,\dots,x_n)\,J(x_1)\cdots J(x_n)\,,
\]
where
\[
W^{(n)}(x_1,\dots,x_n)=\bigl\langle\varphi(x_1)\cdots\varphi(x_n)\bigr\rangle_{\text{conn}}\,.
\]
The functional $W[J]$ is called the generating functional for the
\emph{connected} correlation functions. In the KG theory one can evaluate the
Gaussian functional integral \eqref{eq:ZJ} for $Z[J]$ explicitly and find
\[
W_{\mathrm{KG}}[J]=\frac12\int d^4x\,d^4x'\,J(x)\,D(x-x')\,J(x')\,,
\]
which shows again that in the free theory all connected correlation functions
but the two-point one vanish.

\begin{exercise}
Using $\zeta[J]=\exp W[J]$, check by direct calculation that
\begin{align*}
W^{(2)}(x_1,x_2)&=\bigl\langle\varphi(x_1)\varphi(x_2)\bigr\rangle\,,\\
W^{(4)}(x_1,\dots,x_4)&=\bigl\langle\varphi(x_1)\cdots\varphi(x_4)\bigr\rangle
-\bigl\langle\varphi(x_1)\varphi(x_2)\bigr\rangle\bigl\langle\varphi(x_3)\varphi(x_4)\bigr\rangle\\
&\qquad-\text{two other disconnected terms.}
\end{align*}
\end{exercise}

In many cases it is useful to decompose the connected correlation functions
further into more elementary ``blocks''. Consider the diagrams for the connected
four-point function
\[
\bigl\langle\varphi(x_1)\cdots\varphi(x_4)\bigr\rangle_c=\figfile{fourptconn}
\]
It is clear that some of these diagrams (like (b), (c), and (e)) represent the
corrections to the ``external legs''. This shows that the connected correlation
functions can be expressed through so-called ``amputated'' correlation
functions,
\[
W^{(n)}(x_1,\dots,x_n)=\int\Bigl[\prod_{i=1}^nd^4y_i\,W(x_i-y_i)\Bigr]\,W^{(n)}_{\text{amp}}(y_1,\dots,y_n)
\]
where I used the shorthand
\[
W(x-y)\equiv W^{(2)}(x,y)=\bigl\langle\varphi(x)\varphi(y)\bigr\rangle\,.
\]
Let us give a graphical representation of this relation. Denoting the connected
$n$-point correlation function by an empty blob with $n$ legs,
\[
\bigl\langle\varphi(x_1)\cdots\varphi(x_n)\bigr\rangle_{\text{conn}}=
\diac{\pic{legblob={0.8}{0,45,90,135}{1.1}{fill=white}};
  \draw[densely dotted] (165:1.3) arc[start angle=165,end angle=330,radius=1.3];
  \node at (300:1.75) {\small$n$};}\,,
\]
we have
\[
\diac{\pic{legblob={0.8}{0,45,90,135}{1.1}{fill=white}};
  \draw[densely dotted] (165:1.3) arc[start angle=165,end angle=330,radius=1.3];
  \node at (300:1.75) {\small$n$};}
=\diac{\foreach \a in {0,45,90,135} {\draw[prop] (\a:0.9) -- (\a:2.5) node[vtx] {}; \draw[prop,fill=white] (\a:1.7) circle (0.3);}
  \draw[prop,fill=white] (0,0) circle (0.9); \node at (0,0) {\small amp};
  \draw[densely dotted] (165:1.4) arc[start angle=165,end angle=330,radius=1.4];
  \node at (300:1.85) {\small$n$};}\;;\qquad\qquad
x\ \dia{\draw[prop] (-1.6,0) node[vtx]{} -- (1.6,0) node[vtx]{}; \draw[prop,fill=white] (0,0) circle (0.4);}\ x'=W^{(2)}(x-x')\,.
\]
According to this definition the amputated correlation functions contain
neither external legs nor the corrections to them. Note that by this definition
$W^{(2)}_{\text{amp}}(x,y)=S(x-y)$ is inverse (in the operator sense) to the
two-point function, i.e.
\[
\int d^4y\,S(x-y)\,W(y-x')=\delta^{(4)}(x-x')\,.
\]
In the free theory this is the inverse of the function $D$ i.e.\ just the
(Euclidean) Klein-Gordon operator
\[
S_0(x-y)=\bigl(m^2-\Delta_x\bigr)\,\delta^{(4)}(x-y)\,.
\]

\paragraph{Proper vertices:} Finally, let us define the so called \emph{proper
vertices} (also referred to as the ``one-particle irreducible correlation
functions''). For $n>2$ the $n$-point proper vertex
$-\Gamma^{(n)}(y_1,\dots,y_n)$ is the sum of all diagrams contributing to
$W^{(n)}_{\text{amp}}(y_1,\dots,y_n)$ which cannot be made disconnected by
cutting just one line (please pay attention to the minus sign in this
definition). In particular, in the $\varphi^4$ theory
$\Gamma^{(4)}=-W^{(4)}_{\text{amp}}$ (why?), i.e.\ diagrammatically
\[
\diac{\pic{legblob={0.8}{45,135,225,315}{1.4}{fill=white}};}
=-\ \diac{\pic{dressed={0.8}{45,135,225,315}};}\,,
\]
but for the 6-point functions we have in graphic representation
\[
\diac[0.42]{\pic{legblob={0.8}{0,60,120,180,240,300}{1.4}{fill=white}};}
=-\ \diac[0.42]{\pic{dressed={0.8}{0,60,120,180,240,300}};}
+\Bigl[\ \diac[0.42]{\pic{dressed={0.8}{60,120,180,240,300}};
  \draw[prop] (0.8,0) -- (2.4,0); \draw[prop,fill=white] (1.6,0) circle (0.3);
  \begin{scope}[xshift=3.2cm]\pic{dressed={0.8}{0,60,120,240,300}};\end{scope}
  \node[vtx,minimum size=4pt] at (0.8,0) {}; \node[vtx,minimum size=4pt] at (2.4,0) {};}
\ +\ 9\ \text{similar terms}\Bigr]\,,
\]
where each shadowed blob with $n$ fat points on it stands for the proper vertex
$\Gamma^{(n)}$.

For the reasons that will become clear later the two-point proper vertex
deserves special definition, namely
\[
\Gamma^{(2)}(x,x')=S(x-x')\,.
\]
Its relation to the one-particle irreducible diagrams is as follows. Write
\[
S(x-x')=S_0(x-x')+\Sigma(x-x')\,.
\]
The function $\Sigma(x-x')$ is often referred to as the mass operator. It is
the function $-\Sigma$ (again note the sign) which is equal to the sum of all
one-particle irreducible diagrams for the two-point correlation function.
Indeed, denoting $-\Sigma$ by the two-point shadowed blob we have
\[
x\ \dia{\draw[prop,pattern=north east lines] (0,0) circle (0.6); \node[vtx,minimum size=4pt] at (-0.6,0) {}; \node[vtx,minimum size=4pt] at (0.6,0) {};}\ x'
=\underset{x,x'}{\dia{\draw[prop] (0,0.45) circle (0.45); \node[vtx] at (0,0) {};}}
+\underset{x,x'}{\dia{\draw[prop] (0,0.45) circle (0.45); \draw[prop] (0,1.35) circle (0.45); \node[vtx] at (0,0) {}; \node[vtx] at (0,0.9) {};}}
+\ x\ \dia{\pic{melon};}\ x'+\cdots\,.
\]
The full set of diagrams contributing to $W$,
\[
W(x-x')=x\,\dia{\pic{line};}\,x'
+x\,\dia{\pic{tadline};}\,x'
+x\,\dia{\draw[prop] (-2.2,0) -- (2.2,0); \draw[prop] (-0.8,0.42) circle (0.42); \draw[prop] (0.8,0.42) circle (0.42);
  \foreach \p in {-2.2,-0.8,0.8,2.2} \node[vtx] at (\p,0) {};}\,x'+\cdots\,,
\]
can be obtained by repeating the one-particle irreducible diagrams in a
chainlike manner,
\[
\dia{\draw[prop] (-1.6,0) node[vtx]{} -- (1.6,0) node[vtx]{}; \draw[prop,fill=white] (0,0) circle (0.5); \node[above] at (0,0.5) {\small$W$};}
=\dia{\pic{line};}
+\dia{\draw[prop] (-1.8,0) node[vtx]{} -- (1.8,0) node[vtx]{}; \draw[prop,pattern=north east lines] (0,0) circle (0.5);
  \node[vtx,minimum size=4pt] at (-0.5,0) {}; \node[vtx,minimum size=4pt] at (0.5,0) {};}
+\dia{\draw[prop] (-2.4,0) node[vtx]{} -- (2.4,0) node[vtx]{};
  \foreach \c in {-0.9,0.9} {\draw[prop,pattern=north east lines] (\c,0) circle (0.45);
  \node[vtx,minimum size=4pt] at (\c-0.45,0) {}; \node[vtx,minimum size=4pt] at (\c+0.45,0) {};}}
+\cdots\,.
\]
Explicitly,
\[
W=D-D*\Sigma*D+D*\Sigma*D*\Sigma*D-\cdots\,.
\]
Here $A*B$ denotes a convolution
\[
(A*B)(x-x')=\int d^4y\,A(x-y)\,B(y-x')\,.
\]
Using
\[
S_0*D=I
\]
where $I=\delta^{(4)}(x-y)$ is the unit operator, it is straightforward to check
that $S_0+\Sigma$ is indeed an inverse of $W$.

Similarly to the above generating functionals $\zeta$ and $W$ we can define a
generating functional of the proper vertices. This functional, denoted
$\Gamma[\phi]$, is called the ``Effective Action''.
